% LSPV-1 REPORT::NAEPC2D1::2026-10-19T05:12:00+09:00
% !TEX program = xelatex
\documentclass[11pt]{article}
\usepackage[T1]{fontenc}
\usepackage{lmodern,amsmath,amssymb,amsthm,booktabs,tabularx,xurl,hyperref}
\usepackage[margin=.88in]{geometry}
\newtheorem{theorem}{Theorem}
\newtheorem{lemma}[theorem]{Lemma}
\newtheorem{corollary}[theorem]{Corollary}
\newcommand{\PC}{\ensuremath{\mathsf{PC}}}
\newcommand{\PCR}{\ensuremath{\mathsf{PCR}}}
\newcommand{\EXACT}{\textbf{PROVED/EXACT}}
\newcommand{\KNOWN}{\textbf{KNOWN THEOREM}}
\newcommand{\COMP}{\textbf{COMPUTATIONAL}}
\newcommand{\CAND}{\textbf{CANDIDATE}}
\newcommand{\OPEN}{\textbf{OPEN}}
\newcommand{\RET}{\textbf{FALSIFIED/RETIRED}}
\setlength{\emergencystretch}{3em}
\hypersetup{hidelinks}
\title{NAEPC2D1:\\A Sharp Degree-Two Retraction for the Pin-Free NAE Gauge Compiler}
\author{Tetrahedral Proof Transport / LSPV-1}
\date{19 October 2026}
\begin{document}
\maketitle

\begin{center}\fbox{\parbox{.94\linewidth}{
\textbf{Status.}  For the exact pin-free gauge-chain compiler of
NAEGAUGECHAIN1, every degree-$D$ \PC{} or \PCR{} refutation over every field
projects to a source refutation of degree at most $2D$.  This sharpens the
prior brief's $2D+4$ ledger and yields an all-field linear-degree,
exponential-expanded-size proper signed-NAE family from the verified
graph-colouring source.  The theorem is \EXACT{} at the stated proof-system
model, pending independent review; the 364-check audit is \COMP.  The naive
linewise monomial-size projection can expand by exactly $4^D$ and is
\RET{} as a linear-size transfer.  No accessible source defines the FTNC
ports, ONE3/NOT-ONE3 synchronization, or REVFIX2, so factors $7$, $11$, and
the additive loss $6$ remain \OPEN.  P versus NP remains \OPEN.
}}\end{center}

\section{PREDECESSOR WRITE-RECOVERY PREFLIGHT}
The newest synthesis item is
\path{research/main/2026-10-19_0450_MAIN_RELATIVEGAUGEFOLD23_BRIEF.md}.
It is a substantive Markdown brief, but it is not a canonical MAIN \texttt{.tex}
artifact and has neither the required MAIN header nor final MAIN sentinel.
The last verified canonical MAIN located is
\path{research/main/2026-10-19_0155_MAIN_NAELIFTKW2.tex}.
Classification: \textbf{MALFORMED/NONCANONICAL PREDECESSOR WRITE}.  LSPV-5
should retry the exact full 04:50 body if it exists; this run does not
reconstruct that body from the brief and did not block on recovery.

\section{CORPUS REVIEWED}
The live GitHub tree and current root were surveyed.  The run read
\path{AGENTS.md}, \path{CODEX.md}, the full PNP-research skill and protocols,
all task specifications, the TETRAFOLD/REVFIX directive, NAELIFTKW2,
FTNCNAESLICE1, NAEGAUGEHIGHGIRTH1, TETRAFOLDGAP1, TETRABOUNDARY68,
TETRAPINORBIT1, NAEGAUGECHAIN1, the 00:35 PC22\_ALLFIELD1 report,
PCRETRACT1, and CP22LIFT1.  Repository searches covered PC22\_TRANSFER1,
REVFIX2, ONE3/NOT-ONE3, the claimed degree factors $7/11$, and additive loss
$6$.  Only an attribution inside code was available for PC22\_TRANSFER1; no
REVFIX2 or ONE3/NOT-ONE3 definition was present.  Those branches were stopped
rather than replaced by schematic gates.

Root \path{main.tex}, blob
\texttt{783a11343586a4380ec8559be1ac48218fb322ea}, contains a literal
``Full content continues'' placeholder and no LSPV-1 ownership markers.
AGENTS.md therefore forbids specialist reconciliation of that file.

\section{CLAIM CARD}
\begin{description}
\item[Claim ID.] NAE-PC-DEGREE-RETRACTION-2D1.
\item[Status before run.] \CAND{}; the 04:50 brief claimed $2D+4$ without a
canonical standalone proof artifact.
\item[Quantifiers.] For every field $\mathbb K$, every unsatisfiable nonempty
proper signed 3-CNF $F$, every standard \PC$_\mathbb K$ or
\PCR$_\mathbb K$ refutation $\pi$ of the exact NAEGAUGECHAIN1 target $T(F)$,
if $\deg\pi=D$, then $F$ has a refutation of degree at most $2D$.
\item[Model.] Explicit expanded-polynomial \PC/\PCR, Boolean axioms included;
NAE has the two degree-three forbidden-row axioms.  Degree is the maximum
total degree of a displayed polynomial.
\item[SAME\_REALIZATION.] The target formula, target proof, substitution,
and projected source proof are one explicit algebraic transformation.  No
cost is attributed to an arbitrary SAT execution.
\item[Falsifier.] A target axiom whose image is outside the source Boolean
ideal; a multiplication step needing degree above $2D$; a characteristic-two
failure; or an uncharged assertion of polynomial monomial size.
\end{description}

\section{IDEAS PROPOSED BEFORE LITERATURE}
Four distinct mechanisms were fixed before external searching.
\begin{enumerate}
\item \textbf{Degree-two algebraic retraction.}  Set every gauge and equality
auxiliary to zero and substitute
$u_i=(1-a_i)(1-b_i)$, then simulate every proof inference.
\item \textbf{Gauge restriction/proof by cases.}  Restrict the common gauge
to zero first, and use a second gauge branch only if the direct restriction
does not map all axioms into the source ideal.  Falsifier: branch recombination
cost exceeds a constant additive degree.
\item \textbf{Direct target degree-size route.}  Transfer only degree, then
apply the size-degree theorem to the target itself, avoiding any claim that
expanded source proofs are short.
\item \textbf{Full FTNC/REVFIX local inverse.}  Search for a bounded-degree
substitution through signed FTNC ports, ONE3/NOT-ONE3, and circular feedback.
This branch is valid only with the exact project axioms; missing definitions
are a stopping condition.
\end{enumerate}

\section{PRIMARY LITERATURE CHECK}
Conneryd, de Rezende, Nordstr\"om, Pang and Risse prove linear \PC{} degree
for non-3-colourability of sparse random regular and Erd\H{o}s--R\'enyi
graphs over any field \cite{Conneryd}.  The project source adapter and the
field-independent graph event are inherited exactly from PC22\_ALLFIELD1;
this report does not re-prove their graph theorem.

Mik\v{s}a--Nordstr\"om give clustered-expansion degree lower bounds for
\PC/\PCR{} and explicitly use the Impagliazzo--Pudl\'ak--Sgall connection from
degree to size \cite{MN,IPS}.  The transferable statement here is only the
bounded-degree polynomial-system size-degree theorem.  It says nothing about
arbitrary Boolean circuits or deterministic SAT time.

Low-degree substitution is an exact algebraic proof-system mechanism in this
run.  Guarded restrictions, interpolation, and extension-variable lower bounds
do not automatically apply to the absent REVFIX2 system.  Classification:
the graph-colouring theorem and size-degree theorem are \textbf{THEOREM
TRANSFERS} at their exact models; PC-size-to-general-runtime and
gate-encoding-to-hardness are \textbf{FAILED TRANSFERS}.  The post-quantum
quadratic-relation prompt points in the opposite direction here: the low-rank
quadratic relation is exploitable by the proof projection, so it is an
algorithmic weakness, not a hardness source.

\section{EXACT GAUGE-CHAIN SYSTEM}
Write
\[
F=\bigwedge_{i=1}^m(a_i\vee b_i\vee c_i),
\]
where $a_i,b_i,c_i$ are literal truth polynomials.  NAEGAUGECHAIN1 defines
\begin{align*}
A_i&=\operatorname{NAE}(a_i,b_i,u_i),\\
B_i&=\operatorname{NAE}(\neg u_i,c_i,z_i),\\
E_i&=\operatorname{NAE}(z_i,\neg z_{i+1},e_i)\wedge
\operatorname{NAE}(z_i,\neg z_{i+1},\neg e_i).
\end{align*}
For literal-value polynomials $r,s,t$, the NAE axiom pair is
\[
rst=0,\qquad (1-r)(1-s)(1-t)=0.                 \tag{1}
\]
The source clause axiom is
\[
C_i=(1-a_i)(1-b_i)(1-c_i)=0.                    \tag{2}
\]

Define a substitution $\sigma$ by
\[
\sigma(x)=x\quad\text{for source variables},\qquad
\sigma(z_i)=\sigma(e_i)=0,
\qquad
\sigma(u_i)=(1-a_i)(1-b_i).                     \tag{3}
\]

\begin{lemma}[Axiom images over every field]\label{lem:axioms}
Modulo the source Boolean axioms, every target axiom maps under $\sigma$ to
zero or to one source clause axiom.  Every image has a source derivation of
degree at most four.
\end{lemma}
\begin{proof}
Put $p=(1-a)(1-b)$.  The two axioms for $\operatorname{NAE}(a,b,u)$
map to
\[
abp=ab(1-a)(1-b)=0,
\qquad
(1-a)(1-b)(1-p)=p(1-p)=0.                       \tag{4}
\]
For $\operatorname{NAE}(\neg u,c,z)$ with $z=0$, the all-one axiom
maps to zero and the all-zero axiom maps exactly to
\[
p(1-c)=(1-a)(1-b)(1-c)=C_i.                     \tag{5}
\]
Every equality-chain NAE clause maps to the already-satisfying literal rows
$(0,1,0)$ or $(0,1,1)$, so both polynomials in (1) are zero.

It remains to derive the target Boolean axiom for $u$.  Let
$r=1-a$, $s=1-b$, so $p=rs$.  In the ordinary polynomial ring,
\[
p^2-p=(r^2-r)s^2+r(s^2-s).                       \tag{6}
\]
For either polarity of a source literal, $r^2-r$ is the corresponding
underlying variable Boolean axiom; likewise for $s$.  Multiplication by a
literal is simulated by a variable multiplication and a linear combination.
Thus (6) has a degree-four \PC{} derivation.  The identity is integral and
therefore survives characteristic two.  Constants and source-variable Boolean
axioms are immediate.  In \PCR{}, complement rails are mapped consistently
to $1-\sigma(y)$; the complementarity and product axioms reduce to the same
Boolean consequences.
\end{proof}

\begin{theorem}[Sharp $2D$ proof projection]\label{thm:2d}
For every field $\mathbb K$ and every degree-$D$ \PC$_\mathbb K$ or
\PCR$_\mathbb K$ refutation of $T(F)$, there is a refutation of $F$ of degree
at most $2D$.  Consequently
\[
\deg_{\PC_\mathbb K} T(F)\ge
\tfrac12\deg_{\PC_\mathbb K}F,
\qquad
\deg_{\PCR_\mathbb K} T(F)\ge
\tfrac12\deg_{\PCR_\mathbb K}F.               \tag{7}
\]
\end{theorem}
\begin{proof}
Apply $\sigma$ to every proof line.  Linear combinations commute with
$\sigma$.  If a target inference multiplies a line $P$ by variable $y$, then
$\sigma(y)$ is a constant, a source variable, a source literal, or a product
of two source literals.  From $\sigma(P)$, each such product is obtained using
at most two variable-multiplication steps plus linear combinations.  A target
line of degree $d$ therefore expands only into source polynomials of degree at
most $2d$; intermediate simulation lines obey the same bound.  Lemma
\ref{lem:axioms} initializes all axiom lines in degree at most four.  Since a
refutation using the degree-three NAE axioms has $D\ge3$, degree four is at
most $2D$.  The final constant $1$ is fixed by $\sigma$.  The \PCR{} argument
uses the same simulation for both rails.
\end{proof}

This removes the additive four from the previous brief: it was a safe but
unnecessary global allowance for the local Boolean derivation.

\section{MONOMIAL-SIZE FALSIFIER AND ALL-FIELD CONSEQUENCE}
\begin{lemma}[Exact expansion firewall]
For pairwise disjoint source variables, a target monomial containing $d$
distinct splitter variables has an image with exactly $4^d$ multilinear
monomials and degree $2d$ over every field.
\end{lemma}
\begin{proof}
The image is
\[
\prod_{j=1}^d(1-x_j)(1-y_j).
\]
Choosing independently either $1$ or the variable term from every linear
factor gives $4^d$ distinct supports.  Every coefficient is $\pm1$, hence
nonzero in every field, including characteristic two.
\end{proof}

Therefore the projected expanded size of a degree-$D$ line with $S$ monomials
has only the upper bound $S4^D$ before cancellations.  A linear-size proof
projection is not established.  This does not weaken Theorem~\ref{thm:2d},
because a degree lower bound needs only existence of the projected proof.

For the sparse $6$-regular graph family in PC22\_ALLFIELD1, let $P(F_G)$ be
the verified proper-3CNF adapter.  It has $n_P=15h$ variables and $m_P=25h$
clauses.  NAEGAUGECHAIN1 therefore gives
\[
N=90h-1,\qquad M=100h-2.                         \tag{8}
\]
Conneryd et al. plus the verified degree-three colouring adapter imply
$\deg P(F_G)=\Omega(h)$ over every field on the stated asymptotically
almost-sure non-3-colourable event.  Theorem~\ref{thm:2d} gives
$\deg T(P(F_G))=\Omega(h)=\Omega(N)$.  The standard size-degree theorem then
gives explicit monomial size $2^{\Omega(N)}$, equivalently
$2^{\Omega(L/\log L)}$ for indexed bit length $L=\Theta(N\log N)$.
This is a restricted proof-system corollary, not a deterministic SAT-time
lower bound.

\section{TESTS AND COUNTEREXAMPLES}
The dependency-free checker
\path{experiments/review/nae_pc_degree_retraction_audit.py} uses exact sparse
polynomials in the Boolean quotient.  It exhausts all eight source-literal
polarity patterns, all 64 signed source truth rows, equality-gadget images,
splitter Boolean axioms, and fields $2,3,5,7,11$.  It also computes disjoint
splitter expansion through $d=8$.  Result:
\[
364\text{ assertions},\qquad 0\text{ failures},\qquad 0.616\text{ seconds}.
\]
At $d=8$, every field has exactly $65{,}536=4^8$ monomials and degree $16$.
There is no random seed.  Script SHA-256 is
\texttt{2939c1652f3564951b451b2267f30de70c93ab014bb5f5a2e986d53fb30f1ac4};
raw JSON SHA-256 is
\texttt{7966ae9939699de2237a9eefcae8f0ea05b8f3a7edc731748033317e52beeae6}.
The finite audit is \COMP{} and supports, but does not replace, the symbolic
proof.

Hostile controls were literal renaming, all polarity switches,
characteristic two, constants, equality auxiliaries, and compact
representation.  The compact product representation is the decisive warning:
expanded monomial size can be exponential while an arithmetic circuit for the
same polynomial is linear.

\section{DUPLICATION / CONTRADICTION AUDIT}
The 04:50 brief already states the degree-two substitution and a $2D+4$ bound;
that is adopted as prior work.  The new contribution is the exact axiom-by-
axiom proof, removal of the additive four, independent all-field audit, exact
$4^D$ size firewall, and the compact all-field graph-colouring consequence
with counts (8).  PC22\_ALLFIELD1's variable/zero projection is a different
degree-one compiler and retains its stronger $O(S)$ monomial-size transport.
PCRETRACT1 concerns a different exact-$(2,2)$ Tseitin family and its
characteristic-two escape.  No contradiction is created.

The claimed factors $7$, $11$, and additive loss $6$ have no accessible
defining proof or exact gate system.  They are not inherited.  The earlier
$n+22m,36m$ local-pin claim is likewise not used.  Exact-$(2,2)$ compiler,
Resolution, CP, and graph-colouring PC/PCR results are preserved at their
actual scopes.

\section{WHAT CHANGED RELATIVE TO ROOT MAIN}
Root currently predates the tetrahedral redirect and is structurally
truncated.  The owner-ready change is: add NAEPC2D1 as a proof-system-only
entry, with sharp factor two, exact exponential expansion warning, and the
all-field proper signed-NAE corollary.  Do not change executive OPEN status.
Because the LSPV-1 markers are absent, no root edit was attempted.

\section{ADOPTED / REJECTED / RETIRED PRIOR COMPONENTS}
\begin{itemize}
\item \textbf{ADOPTED:} NAEGAUGECHAIN1's exact compiler and witness map;
PC22\_ALLFIELD1's graph-colouring adapter and source hypotheses; the standard
\PC/\PCR{} size-degree theorem.
\item \textbf{REJECTED:} silently treating a semantic NAE identity as a formal
ring identity without Boolean-axiom derivations.
\item \textbf{RETIRED:} an $O(S)$ linewise expanded-monomial projection for
the quadratic splitter substitution.
\item \textbf{OPEN:} FTNC, ONE3/NOT-ONE3, and REVFIX2 transport; any
restricted-proof-to-arbitrary-machine implication.
\end{itemize}

\section{STALE OR OVERSTATED CLAIMS FOUND}
The 04:50 brief is noncanonical and its $2D+4$ bound is correct but loose.
Any claim that it closes the full FTNC/REVFIX path, preserves shortest proof
size, or implies general SAT time would be overstated.  Root \texttt{main.tex}
is not a complete living synthesis because of its literal placeholder and
missing specialist markers.

\section{BARRIERS, STATUS, AND NEXT HANDOFF}
This is an algebraic proof-system simulation.  It relativizes only in the
irrelevant sense that no arbitrary computation is modelled; no
nonrelativizing or nonalgebrizing hinge is supplied.  Natural Proofs is not
invoked.  A polynomial-time SAT decider need not output a \PC/\PCR{}
refutation, and a compact arithmetic or Boolean circuit need not expand the
projected monomials.  Thus P versus NP remains \OPEN.

\textbf{Highest-value next obligation.}  Recover the exact signed FTNC port
polynomials and the ONE3/NOT-ONE3 and REVFIX2 axiom lists.  For each target
variable, specify a source polynomial image and verify every gate axiom in the
source ideal over characteristic two and an arbitrary odd characteristic.
Only then recompute, rather than inherit, the proposed factors $7$, $11$, and
additive loss $6$.

\begin{thebibliography}{4}
\bibitem{Conneryd} J. Conneryd, S. F. de Rezende, J. Nordstr\"om, S. Pang,
and K. Risse, ``Graph Colouring Is Hard on Average for Polynomial Calculus
and Nullstellensatz,'' FOCS 2023; full version arXiv:2503.17022.
\url{https://arxiv.org/abs/2503.17022}
\bibitem{MN} M. Mik\v{s}a and J. Nordstr\"om, ``A Generalized Method for
Proving Polynomial Calculus Degree Lower Bounds,'' CCC 2015, LIPIcs 33,
467--487; arXiv:1505.01358.
\url{https://arxiv.org/abs/1505.01358}
\bibitem{IPS} R. Impagliazzo, P. Pudl\'ak, and J. Sgall, ``Lower Bounds for
the Polynomial Calculus and the Gr\"obner Basis Algorithm,'' Computational
Complexity 8 (1999), 127--144; ECCC TR97-042.
\url{https://eccc.weizmann.ac.il/eccc-reports/1997/TR97-042/}
\bibitem{Project} LSPV-4, ``NAEGAUGECHAIN1,'' canonical project report,
19 October 2026.
\end{thebibliography}
\end{document}
% END_OF_VERIFIED_REPORT::LSPV-1::NAEPC2D1
